\subsection{The Hahn-Banach theorem (analytic form)}

THEOREM 1 (Hahn-Banach). --- Let p be a sub-linear function on a vector space E. Let V be a vector subspace of E and f a linear form on V such that, for all \(y \in V\) , we have \(f(y) \leqslant p(y)\) . Then there exists a linear form h on E that is an extension of f and such that \(h(x) \leqslant p(x)\) for \(x \in E\) .

The set of pairs \((x, a)\) such that \(p(x) \leqslant a\) is a convex subset P of the vector space \(E_{1} = E \times R\) (II, p.~17, prop. 19), and it is clearly a pointed cone. Let \(V_{1}\) be the subspace \(V \times R\) of \(E_{1}\) and \(g(y, a) = -f(y) + a\) for each point \((y, a) \in V_{1}\) . Then g is a positive linear form for the preorder structure on \(V_{1}\) defined by \(P \cap V_{1}\) ; for if \((y, a) \in P \cap V_{1}\) , then \(a \geqslant p(y) \geqslant f(y)\) , therefore \(g(y, a) \geqslant 0\) . Next let \((x, a) \in E_{1}\) ; we show that \((x, a)\) is less than a point of \(V_{1}\) for the preorder defined by P. If \((x', a') \in V_{1}\) then \((x, a) \leqslant (x', a')\) if, and only if, \(p(x' - x) \leqslant a' - a\) , taking \(a' \geqslant p(-x) + a\) , we see that \((0, a')\) of \(V_{1}\) satisfies the requirements. Thus we can apply prop. 1 of II, p.~21; there is a linear form u on \(E_{1}\) extending g and positive for the preorder defined by P. Therefore \(u(0, 1) = g(0, 1) = 1\) and u is of the form \(u(x, a) = -h(x) + a\) , where h is a linear form on E that extends f; further, for all \(x \in E\) and all \(a \geqslant p(x)\) , we have \(h(x) \leqslant a\) , therefore \(h(x) \leqslant p(x)\) . Q.E.D.

COROLLARY 1. --- Let p be a semi-norm on the vector space E. Let V be a vector subspace of E and f a linear form on V such that \(|f(y)| \leqslant p(y)\) for all \(y \in V\) . Then there exists a linear form h defined on E which is an extension of f and is such that \(|h(x)| \leqslant p(x)\) for \(x \in E\) .

For a semi-norm \(q\) and a linear form \(g\) on \(\mathbf{E}\) , the relation \(g \leqslant q\) is the same as \(|g| \leqslant q\) . The corollary follows from th. 1.

COROLLARY 2. --- Let p be a semi-norm on the vector space E. Given a point \(x_{0} \in E\) , there exists a linear form f defined over E, such that \(f(x_{0}) = p(x_{0})\) and that \(|f(x)| \leqslant p(x)\) for all \(x \in E\) .

Apply cor. 1 to the vector subspace, V, generated by \(x_0\) and to the linear form \(\xi x_0 \mapsto \xi p(x_0)\) defined over V.

COROLLARY 3. --- Let V be a vector subspace of the normed space E and let f be a continuous linear form over V; then there exists a continuous linear form h defined over E which extends f and is of the same norm (GT, X, § 3.2).

Apply cor. 1, taking \(p(x) = \| f \| \cdot \| x \|\) , which gives \(\| h \| \leqslant \| f \|\) ; but clearly \(\| h \| \geqslant \| f \|\) , and the corollary follows.

The conclusion of cor. 3 is not necessarily valid for continuous linear mappings of a normed space into an arbitrary normed space (IV, p.~55, exerc. 16, c) and V, p.~65, exerc. 22).

